% SPDX-License-Identifier: MIT
% From h0ffmann/closed-source format/tex/template.tex at 6b3aae5; one change: the author line prints the ORCID iD.
% Paper Fabric article template for pandoc. Compiles with XeLaTeX (the local and CI build) and
% with pdfLaTeX, which is the engine arXiv runs on submitted sources (just arxiv). Modelled on the layout of Sandy Maguire's
% book templates (isovector/thinking-with-types, format/tex/template.tex): one hand-written
% template, metadata from the paper's metadata.md, and the Git revision stamped on every page.
% The CSL and table blocks follow pandoc's default LaTeX template (pandoc -D latex).
\documentclass[10pt,a4paper,]{article}

\usepackage{iftex}
\ifPDFTeX
\usepackage[T1]{fontenc}
\usepackage[utf8]{inputenc}
\usepackage{lmodern}
\usepackage{textcomp}
\usepackage{amssymb}
\else
\usepackage{fontspec}
\usepackage{unicode-math}
\defaultfontfeatures{Ligatures=TeX}
\fi
\usepackage[top=14mm,bottom=17mm,left=18mm,right=18mm]{geometry}
\usepackage{microtype}
\usepackage{xcolor}
\definecolor{fabric}{HTML}{C2410C}
\usepackage{graphicx}
\usepackage{calc}
\usepackage{fancyhdr}
\usepackage{titling}
\usepackage{enumitem}
\setlist{itemsep=2pt,topsep=4pt}
\usepackage[brazilian]{babel}
\usepackage{csquotes}
\setlength{\parindent}{0pt}
\setlength{\parskip}{6pt plus 2pt minus 1pt}
\setlength{\emergencystretch}{3em}
\providecommand{\tightlist}{\setlength{\itemsep}{0pt}\setlength{\parskip}{0pt}}


\usepackage{longtable,booktabs,array}
\usepackage{etoolbox}
\makeatletter
\patchcmd\longtable{\par}{\if@noskipsec\mbox{}\fi\par}{}{}
\makeatother
\usepackage{footnotehyper}
\makesavenoteenv{longtable}

% citeproc citations and the reference list, as in pandoc 3.7's common.latex (pandoc >= 3.1.8
% writes \citeproc and \bibitem). The \small in the list is ours.
\NewDocumentCommand\citeproctext{}{}
\NewDocumentCommand\citeproc{mm}{%
  \begingroup\def\citeproctext{#2}\cite{#1}\endgroup}
\makeatletter
 \let\@cite@ofmt\@firstofone
 \def\@biblabel#1{}
 \def\@cite#1#2{{#1\if@tempswa , #2\fi}}
\makeatother
\newlength{\cslhangindent}
\setlength{\cslhangindent}{1.5em}
\newlength{\csllabelwidth}
\setlength{\csllabelwidth}{3em}
\newenvironment{CSLReferences}[2] % #1 hanging-indent, #2 entry-spacing
 {\small\begin{list}{}{%
  \setlength{\itemindent}{0pt}
  \setlength{\leftmargin}{0pt}
  \setlength{\parsep}{0pt}
  \ifodd #1
   \setlength{\leftmargin}{\cslhangindent}
   \setlength{\itemindent}{-1\cslhangindent}
  \fi
  \setlength{\itemsep}{#2\baselineskip}}}
 {\end{list}}
\newcommand{\CSLBlock}[1]{\hfill\break\parbox[t]{\linewidth}{\strut\ignorespaces#1\strut}}
\newcommand{\CSLLeftMargin}[1]{\parbox[t]{\csllabelwidth}{\strut#1\strut}}
\newcommand{\CSLRightInline}[1]{\parbox[t]{\linewidth - \csllabelwidth}{\strut#1\strut}}
\newcommand{\CSLIndent}[1]{\hspace{\cslhangindent}#1}

\setcounter{secnumdepth}{3}

\AtBeginEnvironment{longtable}{\small}
\newfontfamily\markfont{DejaVuSans.ttf}
\catcode`\⚠=\active
\protected\def⚠{{\markfont\char"26A0\relax}}
\makeatletter
\@ifpackageloaded{subcaption}{}{\usepackage{subcaption}}
\@ifpackageloaded{caption}{}{\usepackage{caption}}
\captionsetup[subfigure]{margin=0.5em}
\AtBeginDocument{%
\renewcommand*\figurename{Figura}
\renewcommand*\tablename{Tabela}
}
\AtBeginDocument{%
\renewcommand*\listfigurename{List of Figures}
\renewcommand*\listtablename{List of Tables}
}
\newcounter{pandoccrossref@subfigures@footnote@counter}
\newenvironment{pandoccrossrefsubfigures}{%
\setcounter{pandoccrossref@subfigures@footnote@counter}{0}
\begin{figure}\centering%
\gdef\global@pandoccrossref@subfigures@footnotes{}%
\DeclareRobustCommand{\footnote}[1]{\footnotemark%
\stepcounter{pandoccrossref@subfigures@footnote@counter}%
\ifx\global@pandoccrossref@subfigures@footnotes\empty%
\gdef\global@pandoccrossref@subfigures@footnotes{{##1}}%
\else%
\g@addto@macro\global@pandoccrossref@subfigures@footnotes{, {##1}}%
\fi}}%
{\end{figure}%
\addtocounter{footnote}{-\value{pandoccrossref@subfigures@footnote@counter}}
\@for\f:=\global@pandoccrossref@subfigures@footnotes\do{\stepcounter{footnote}\footnotetext{\f}}%
\gdef\global@pandoccrossref@subfigures@footnotes{}}
\@ifpackageloaded{float}{}{\usepackage{float}}
\floatstyle{ruled}
\@ifundefined{c@chapter}{\newfloat{codelisting}{h}{lop}}{\newfloat{codelisting}{h}{lop}[chapter]}
\floatname{codelisting}{Listing}
\newcommand*\listoflistings{\listof{codelisting}{List of Listings}}
\makeatother

\usepackage{xurl}
\usepackage{hyperref}
\hypersetup{
  pdftitle={Modelos de ondas em GPUs},
  pdfauthor={Matheus Hoffmann; Pedro Veras Guimarães; Pedro Cruz},
  pdfkeywords={WAVEWATCH III, reprodutibilidade bit a
bit, Kokkos, Rust, porte para GPU, agentes de programação},
  colorlinks=true, linkcolor=fabric, citecolor=fabric, urlcolor=fabric,
  pdfcreator={Paper Fabric (pandoc + XeLaTeX)}}
\urlstyle{same}

% The build writes the short Git hash and commit count here (justfile, build/.current-revision).
\newcommand{\revision}{\IfFileExists{.current-revision}{\input{.current-revision}}{unversioned}}

\pagestyle{fancy}
\fancyhf{}
\renewcommand{\headrulewidth}{0pt}
\fancyfoot[L]{\footnotesize\color{gray}Modelos de ondas em
GPUs \textperiodcentered{} v0 (preprint, sem revisão por pares)}
\fancyfoot[R]{\footnotesize\color{gray}rev.~\revision{} \textperiodcentered{} \thepage}

\begin{document}

\begin{center}
{\LARGE\bfseries Modelos de ondas em GPUs\par}
\vspace{4pt}
{\large Portes do WAVEWATCH III em C++/Kokkos e em Rust que precisam dar
os bits do Fortran\par}
\vspace{12pt}
{\large Matheus
Hoffmann}\,\href{https://orcid.org/0009-0009-1056-7661}{\footnotesize\color{fabric}ORCID 0009-0009-1056-7661}\\
{\small\itshape Escola Politécnica, Universidade Federal do Rio de
Janeiro, Brasil}\par\vspace{4pt}
{\large Pedro Veras
Guimarães}\,\href{https://orcid.org/0000-0002-1686-4722}{\footnotesize\color{fabric}ORCID 0000-0002-1686-4722}\\
{\small\itshape Coorientador do primeiro autor. LabECO, Departamento de
Engenharia Mecânica, Universidade Federal de Santa Catarina,
Brasil}\par\vspace{4pt}
{\large Pedro
Cruz}\,\href{https://orcid.org/0000-0002-2019-133X}{\footnotesize\color{fabric}ORCID 0000-0002-2019-133X}\\
{\small\itshape Orientador do primeiro autor. Grupo de Teleinformática e
Automação (GTA), DEL, Escola Politécnica, Universidade Federal do Rio de
Janeiro, Brasil}\par\vspace{4pt}
\vspace{6pt}
{\small 2026-11-11 \textperiodcentered{} v0 (preprint, sem revisão por
pares) \textperiodcentered{} rev.~\revision\par}
\end{center}

\begin{quote}
\small\textbf{Resumo.} Este texto relata dois portes do modelo espectral
de ondas WAVEWATCH III (WW3). Agentes de programação escreveram os dois
portes, e os dois portes têm uma regra de aceitação só: o porte precisa
dar os mesmos bits do Fortran. Uma tolerância é um diagnóstico e nunca é
o veredito. O \emph{ww3-gpu} mantém o WW3, e o Fortran chama código C++
escrito com Kokkos, atrás de uma flag de execução para cada rotina. A
referência do ww3-gpu é o caminho Fortran do mesmo executável. O
\emph{chaceon} substitui o WW3 por Rust, compilado para a CPU, para
SPIR-V (rust-gpu) e para PTX (Rust-CUDA). Nenhum Fortran roda no
processo do chaceon, e a referência do chaceon é um WW3 compilado em
separado. O ww3-gpu porta a aproximação de interação discreta (DIA) da
transferência entre quatro ondas. Na CPU, a DIA portada é idêntica byte
a byte nas 100 entradas seriais e nas 25 entradas OpenMP da matriz de
regressão do WW3 que selecionam a DIA. A DIA ocupa 48,4 \% do tempo
próprio em uma thread, num caso de 9.520 pontos. O chaceon roda
\texttt{ww3\_grid}, \texttt{ww3\_strt}, \texttt{ww3\_prep} e
\texttt{ww3\_shel} sem Fortran, e 96 arquivos de saída são idênticos
byte a byte à referência. Nenhum dos dois portes rodou o modelo inteiro
numa GPU real, e este texto não afirma ganho de velocidade. O texto
compara como cada projeto define e confere a própria evidência.

\textbf{Palavras-chave:} WAVEWATCH III; reprodutibilidade bit a
bit; Kokkos; Rust; porte para GPU; agentes de programação.
\end{quote}

\section*{Prefácio}\label{prefuxe1cio}
\addcontentsline{toc}{section}{Prefácio}

Este projeto existe para apoiar o Ondas \& Ventos
(\citeproc{ref-ondasventos}{Universidade Federal de Santa Catarina
2026}), um sistema operacional de previsão de ondas da Universidade
Federal de Santa Catarina (UFSC). O LabECO da UFSC roda o WW3 de forma
operacional para a ReNOMO, a Rede Nacional de Observação e Monitoramento
Oceânico (\citeproc{ref-renomo}{LabECO/UFSC 2025}). A ReNOMO começou em
2023, com a chamada CNPq/MCTI/Finep 062/2022. Mais de 100 pesquisadores
de 31 instituições brasileiras participam da rede
(\citeproc{ref-renomo}{LabECO/UFSC 2025}).

O sistema precisa rodar o modelo de ondas a cada ciclo de previsão, e o
tempo de execução do WW3 limita o ciclo. O sistema também planeja forçar
o modelo de ondas com mais de um modelo de vento: o IFS do ECMWF, o GFS
do NCEP e o WeatherNext do Google ⚠. Cada modelo de vento acrescenta uma
execução completa do WW3, e por isso um conjunto de modelos de vento
multiplica o custo do ciclo. Um WW3 mais rápido em GPU, com os mesmos
bits, torna essas execuções viáveis sem mudança na previsão. O ww3-gpu e
o chaceon trabalham para chegar a esse WW3, e este texto não afirma
ganho de velocidade.

\setcounter{section}{-1}

\section{Contexto: um modelo oceânico portado por um modelo de
linguagem}\label{sec:context}

Koldunov et al. (\citeproc{ref-koldunov2026}{Koldunov et al. 2026})
relatam o porte de um modelo oceânico inteiro por um modelo de linguagem
grande. O modelo é o FESOM2, um modelo de oceano e gelo marinho, e os
autores trabalham no Alfred Wegener Institute. O Claude Code, com o
Claude Opus 4.7, traduziu o FESOM2 de Fortran para C e depois de C para
C++/Kokkos. As pessoas definiram a estratégia, revisaram os planos e os
diffs e encontraram erros de física e de validação. O agente leu,
traduziu, construiu os ambientes de teste, depurou, mediu o desempenho e
escreveu a documentação.

A etapa em C é o centro do método de Koldunov et al.~O código C limpo é
a referência, e cada kernel Kokkos mantém o original em C como
``gêmeo''. A regra de aceitação muda a cada etapa. Os autores comparam o
modelo em C com o Fortran por estatística: em cinco anos, a diferença
RMS da média é 0,006 °C na temperatura da superfície do mar. Só o modelo
Kokkos Serial é idêntico byte a byte, e só ao modelo em C. Os autores
escrevem que a concordância bit a bit entre Fortran e C ``across
compilers and arithmetic orderings is neither attainable nor
necessary''. Na GPU, os autores aceitam diferenças até um teto
registrado para cada campo.

O sucessor do WW3 também está em desenvolvimento. A NCEP Office Note 525
(\citeproc{ref-on525}{Tolman 2025}) escolheu uma reescrita de baixo para
cima, com um núcleo em C++ primeiro e um núcleo em Rust depois ⚠. O
mesmo relatório rebaixa o Fortran a uma linguagem só de solvers ⚠. O
repositório do WAVEWATCH IV (\citeproc{ref-ww4}{NOAA Environmental
Modeling Center 2026}) está em pré-alfa: em 2026-11-11, o branch
\texttt{develop} tinha 40 commits, de 2025-11-17 a 2026-10-16
(\texttt{875f49a}). Os 97 arquivos-fonte do WAVEWATCH IV são C++ (66
\texttt{.cpp}, 31 \texttt{.h}), e nenhum kernel de física do WW3 está no
repositório ainda.

Por isso, os kernels centrais do modelo de ondas continuam em Fortran
por alguns anos ⚠. Esse é o primeiro motivo para o primeiro autor
continuar a manter o ww3-gpu e o chaceon. Os dois projetos fazem os
kernels Fortran do WW3 rodarem em GPU agora, sem mudança nos bits.

Este texto parte desse resultado e faz uma pergunta mais estreita. Um
agente consegue portar um modelo espectral de ondas e dar os bits do
próprio Fortran? Os dois portes abaixo não têm etapa em C, e a
referência dos dois portes é o Fortran. Os dois portes também usam duas
práticas que o FESOM2 relata: tradução estritamente literal e, para cada
afirmação sobre o Fortran, um \texttt{file:line} com o valor literal
lido ali.

\section{Por que bit a bit}\label{sec:why}

O WW3 integra a equação de balanço de ação

\[\frac{\partial N}{\partial t} + \nabla\cdot(\dot{\mathbf{c}}\,N) = \frac{S}{\sigma},\]

em que o termo-fonte \(S\) é a soma da entrada de vento, da dissipação,
da transferência não linear e dos termos de fundo. Um porte que muda a
ordem de uma soma em ponto flutuante muda \(N\). Um sistema de previsão
validado contra observações precisa então ser validado de novo. Por
isso, os dois projetos mantêm a aritmética do Fortran. Os dois projetos
não reordenam reduções, não usam multiplicação-adição fundida
(\texttt{-ffp-contract=off} nos dois lados) e não alargam tolerâncias.
Uma diferença é defeito do porte até duas execuções do próprio Fortran
darem bytes diferentes no mesmo arquivo.

\section{Duas formas para uma regra}\label{sec:shapes}

\begin{longtable}[]{@{}
  >{\raggedright\arraybackslash}p{(\linewidth - 4\tabcolsep) * \real{0.2200}}
  >{\raggedright\arraybackslash}p{(\linewidth - 4\tabcolsep) * \real{0.4000}}
  >{\raggedright\arraybackslash}p{(\linewidth - 4\tabcolsep) * \real{0.3800}}@{}}
\caption{\label{tbl:shapes}Os dois portes. Fontes: chaceon
\texttt{docs/WW3\_GPU\_PARITY\_202611.md} (2026-11-09) e
\texttt{README.md} em \texttt{0c098fe} (\citeproc{ref-chaceon}{Hoffmann
2026a}); ww3-gpu \texttt{kokkos/PORT\_STATUS.md} em \texttt{58167d7}
(\citeproc{ref-ww3gpu}{Hoffmann 2026c}).}\tabularnewline
\toprule\noalign{}
\begin{minipage}[b]{\linewidth}\raggedright
\end{minipage} & \begin{minipage}[b]{\linewidth}\raggedright
ww3-gpu (C++/Kokkos)
\end{minipage} & \begin{minipage}[b]{\linewidth}\raggedright
chaceon (Rust)
\end{minipage} \\
\midrule\noalign{}
\endfirsthead
\toprule\noalign{}
\begin{minipage}[b]{\linewidth}\raggedright
\end{minipage} & \begin{minipage}[b]{\linewidth}\raggedright
ww3-gpu (C++/Kokkos)
\end{minipage} & \begin{minipage}[b]{\linewidth}\raggedright
chaceon (Rust)
\end{minipage} \\
\midrule\noalign{}
\endhead
\bottomrule\noalign{}
\endlastfoot
Forma & Dentro do WW3: o Fortran chama C++ por um shim \texttt{bind(C)},
uma flag para cada rotina & Substitui o WW3: nenhum Fortran num processo
do chaceon \\
Referência & NOAA-EMC/WW3 \texttt{03adda91}, por um fork que acrescenta
5.687 linhas e não remove nenhuma & MetOcean WW3 7.14,
\texttt{92e2cf66} \\
Paridade na CPU contra & o caminho Fortran do mesmo executável & uma
referência compilada em separado \\
Alvos de dispositivo & Kokkos 5.2.2, backend CUDA (\texttt{sm\_89});
OpenMP na CPU & rust-gpu (SPIR-V, Vulkan) e Rust-CUDA (PTX) \\
Espectro residente no dispositivo & sim, no backend OpenMP & sim, em
SPIR-V no Vulkan de CPU (lavapipe) \\
Tamanho do porte & 62.257 linhas de C++, 355 arquivos (2026-10-14) &
29.229 linhas de porte em Rust, 5.936 linhas de ferramentas \\
Métodos formais & consultas SMT, varreduras exaustivas & uma prova em
Lean 4 sobre a divisão de \texttt{f32} \\
Licença & MIT, kernels em LGPL-3.0 & MPL-2.0 \\
\end{longtable}

\begin{longtable}[]{@{}
  >{\raggedright\arraybackslash}p{(\linewidth - 4\tabcolsep) * \real{0.3053}}
  >{\raggedright\arraybackslash}p{(\linewidth - 4\tabcolsep) * \real{0.3474}}
  >{\raggedright\arraybackslash}p{(\linewidth - 4\tabcolsep) * \real{0.3474}}@{}}
\caption{\label{tbl:repos}Os dois repositórios. Fonte: a API REST do
GitHub (\texttt{repos/h0ffmann/\textless{}repo\textgreater{}} para
\texttt{created\_at} e \texttt{pushed\_at}, \texttt{commits?sha=main}
para os commits), lida em 2026-11-11 às 03h35 UTC. O primeiro commit do
ww3-gpu é mais antigo que o repositório porque o commit foi feito antes
do primeiro push.}\tabularnewline
\toprule\noalign{}
\begin{minipage}[b]{\linewidth}\raggedright
Métrica
\end{minipage} & \begin{minipage}[b]{\linewidth}\raggedright
ww3-gpu
\end{minipage} & \begin{minipage}[b]{\linewidth}\raggedright
chaceon
\end{minipage} \\
\midrule\noalign{}
\endfirsthead
\toprule\noalign{}
\begin{minipage}[b]{\linewidth}\raggedright
Métrica
\end{minipage} & \begin{minipage}[b]{\linewidth}\raggedright
ww3-gpu
\end{minipage} & \begin{minipage}[b]{\linewidth}\raggedright
chaceon
\end{minipage} \\
\midrule\noalign{}
\endhead
\bottomrule\noalign{}
\endlastfoot
Repositório criado & 2026-09-12, 01h31 UTC & 2026-10-23, 03h58 UTC \\
Primeiro commit no \texttt{main} & 2026-09-11, 13h47 UTC,
\texttt{b43a6d3} & 2026-10-23, 03h58 UTC, \texttt{3bad31a} \\
Último commit no \texttt{main} & 2026-11-11, 03h16 UTC, \texttt{efc0211}
& 2026-11-09, 08h53 UTC, \texttt{0c098fe} \\
Último push, qualquer branch & 2026-11-11, 03h31 UTC & 2026-11-09, 08h53
UTC \\
Commits no \texttt{main} & 800 & 110 \\
\end{longtable}

A rota Kokkos mantém todos os programas do WW3 e substitui uma rotina de
cada vez. Por isso, um regtest consegue isolar o efeito de uma rotina. A
rota Rust precisa reconstruir os programas, os formatos de arquivo e o
laço de tempo antes de comparar uma execução inteira. Depois disso,
nenhum Fortran sobra para dividir estado com o porte.

O registro de progresso do ww3-gpu para a NOAA
(\citeproc{ref-noaarecord}{Hoffmann 2026e}) guarda os resultados de CPU
do runtime ww3-gpu inteiro, com o espectro residente, e os limites
desses resultados. O porte da DIA também está publicado sozinho, como
ww3-dia-gpu-v2 (\citeproc{ref-diagpuv2}{Hoffmann 2026d}). O repositório
ww3-dia-gpu-v2 compara o porte com um kernel CUDA anterior da mesma
rotina.

\section{Escadas de evidência}\label{sec:ladders}

Cada projeto dá nome aos próprios degraus. Um degrau num registro não é
a mesma evidência que um degrau no outro registro.

\begin{longtable}[]{@{}
  >{\raggedright\arraybackslash}p{(\linewidth - 4\tabcolsep) * \real{0.2400}}
  >{\raggedright\arraybackslash}p{(\linewidth - 4\tabcolsep) * \real{0.3800}}
  >{\raggedright\arraybackslash}p{(\linewidth - 4\tabcolsep) * \real{0.3800}}@{}}
\caption{\label{tbl:ladders}Degraus de evidência. As contagens do
ww3-gpu vêm de \texttt{scripts/port\_score.py\ -\/-write} (perfil
\texttt{bench\_small}, 120 × 80 pontos, 12 h, uma thread, gprof,
2026-10-14); as contagens do chaceon vêm de
\texttt{just\ port-summary}.}\tabularnewline
\toprule\noalign{}
\begin{minipage}[b]{\linewidth}\raggedright
Degrau
\end{minipage} & \begin{minipage}[b]{\linewidth}\raggedright
ww3-gpu
\end{minipage} & \begin{minipage}[b]{\linewidth}\raggedright
chaceon
\end{minipage} \\
\midrule\noalign{}
\endfirsthead
\toprule\noalign{}
\begin{minipage}[b]{\linewidth}\raggedright
Degrau
\end{minipage} & \begin{minipage}[b]{\linewidth}\raggedright
ww3-gpu
\end{minipage} & \begin{minipage}[b]{\linewidth}\raggedright
chaceon
\end{minipage} \\
\midrule\noalign{}
\endhead
\bottomrule\noalign{}
\endlastfoot
O código existe & \emph{translated}: 316 rotinas, 26,2 \% das linhas de
Fortran & \emph{translated}: o Rust existe, sem cobertura completa \\
Uma execução chama o código & \emph{ported}: 72 rotinas, 99,3 \% do
tempo de execução medido & \emph{routine}: toda chamada capturada é
idêntica bit a bit no host e em SPIR-V \\
Execuções inteiras idênticas, CPU & \emph{regtested}: 14 rotinas, 48,6
\% do tempo de execução & \emph{host}: execuções inteiras idênticas por
\texttt{cmp}; 10.080 linhas \\
Idêntico por uma API de GPU & nenhum & \emph{spirv}: os mesmos arquivos
a partir de SPIR-V no Vulkan de CPU; 5.189 linhas \\
Idêntico numa GPU real & \emph{gpu}: nenhum & \emph{cuda}: nenhum \\
\end{longtable}

\section{O que cada projeto encontrou}\label{sec:found}

\textbf{ww3-gpu: a referência é o código otimizado.} A identidade byte a
byte da DIA precisou de um fato que o texto-fonte não dá. O gfortran
expande \texttt{x**11} como uma cadeia de multiplicações em \texttt{-O0}
e como outra cadeia em \texttt{-O3}. As duas cadeias discordam em
643.427 de 1.000.000 entradas de precisão simples. Antes dessa
descoberta, o diário de pesquisa explicava as diferenças restantes como
arredondamento amplificado pelo WW3. Essa explicação estava errada. Por
isso, a referência do porte é o código que o build do WW3 produz, lido
como GIMPLE ou desmontagem, e não o texto Fortran.

O porte do ww3-gpu não tem testes unitários próprios. Só os testes de
regressão do WW3, comparados pelo \texttt{test.comp}, julgam o porte.

\textbf{chaceon: uma execução sem Fortran.} Os programas Rust
\texttt{ww3\_grid}, \texttt{ww3\_strt}, \texttt{ww3\_prep} e
\texttt{ww3\_shel} rodam em sequência a partir das namelists. Os
programas escrevem o \texttt{mod\_def.ww3}, a forçante, as saídas em
grade e em pontos, o arquivo de trajetória e os dois restarts. Os 96
arquivos são idênticos byte a byte à referência. As fronteiras dão 214
comparações idênticas em 34 casos, e os pós-processadores de trajetória
dão 74. O \texttt{ww3\_shel\ -\/-threads\ N} e um build MPI dão os
mesmos bytes.

Os arquivos continuam idênticos quando o \texttt{W3SRCE}, o
\texttt{W3XYP3} e o \texttt{W3KTP3} são compilados para SPIR-V e rodam
num driver Vulkan de CPU. A enumeração não alcança os \(2^{64}\) pares
de operandos \texttt{f32}. Nesse caso, uma prova em Lean enuncia o fato
sobre a divisão que o porte usa. A prova não diz nada sobre a igualdade
entre Rust e Fortran, porque o Fortran não tem semântica formal
utilizável.

\textbf{Por que o ww3-gpu ainda não tem prova, e o que vem depois.}
Nesta primeira fase, o ww3-gpu não adotou assistente de prova
(\citeproc{ref-adr0001}{Hoffmann 2026f}). Há quatro motivos:

\begin{enumerate}
\def\labelenumi{\arabic{enumi}.}
\tightlist
\item
  Nenhum assistente de prova lê Fortran ou C++. Uma prova trata de um
  modelo que uma pessoa escreve. A prova não diz nada sobre os binários
  do gfortran e do nvcc enquanto ninguém mostrar que o modelo concorda
  com esses binários, e essa concordância é a questão em aberto.
\item
  A igualdade de bits é uma propriedade de dois binários, compilados com
  flags dadas. A igualdade de bits não é uma propriedade de dois
  textos-fonte.
\item
  Para código IEEE-754 sem desvios, a igualdade em todas as entradas é
  decidível. Consultas SMT na teoria \texttt{QF\_FP} e varreduras sobre
  as \(2^{32}\) entradas \texttt{float32} respondem a pergunta em
  segundos ou minutos. Uma prova interativa leva meses para cada rotina.
\item
  Agentes precisam escrever cada prova, e uma pessoa precisa revisar
  cada prova. A prova é um segundo artefato para cada rotina.
\end{enumerate}

Sem prova, o porte da DIA tem a evidência mais forte abaixo de uma
prova. O porte é idêntico byte a byte nas 125 entradas da matriz de
regressão do WW3 que selecionam a DIA (100 seriais, 25 OpenMP). A seção
1 do kernel também passou por uma varredura das \(2^{32}\) entradas
\texttt{float32}, com 0 diferenças (\citeproc{ref-adr0001}{Hoffmann
2026f}). Por isso, os testes do próprio modelo validam o porte da DIA em
toda entrada que chama o porte. Essa validação não é uma prova formal, e
este texto não chama a validação de prova.

Uma primeira camada formal já existe ao lado dos portes. O pacote Lean 4
\texttt{ocean-waves} (\citeproc{ref-oceanwaves}{Hoffmann 2026b}) enuncia
a equação de balanço de ação discretizada e prova estas propriedades,
sem \texttt{sorry} e sem \texttt{axiom}:

\begin{itemize}
\tightlist
\item
  \(\sigma^2(k)\) cresce estritamente com \(k\);
\item
  \(c/2 \le c_g \le c\);
\item
  a condição CFL de profundidade mantém todo número de Courant em
  \([0, 1]\);
\item
  a propagação upwind e uma troca de quadrupletos no formato da DIA
  conservam a ação;
\item
  o passo semi-implícito do termo-fonte mantém a ação positiva para todo
  passo de tempo.
\end{itemize}

O pacote usa aritmética \texttt{Rat} exata num modelo reduzido: uma
dimensão espacial, profundidade fixa e um modelo de Padé para \(\tanh\).
Por isso, o pacote também não diz nada sobre bits. O Agda é o próximo
passo. O Agda não tem tipo de ponto flutuante de 32 bits, então o Agda
não consegue provar um fato sobre bits. O papel do Agda é uma semântica
de referência total e executável de cada termo-fonte. Cada índice de
tabela terá um tipo \texttt{Fin}, e toda leitura fica dentro dos limites
para toda grade, não só para uma grade. As varreduras, as consultas ao
solver e os testes de regressão do WW3 continuam a decidir os bits.

\section{Desafios: os testes de regressão são
grandes}\label{sec:challenges}

Os testes de regressão do WW3 são o único juiz do porte Kokkos, e esses
testes são grandes. A Tabela~\ref{tbl:regtests} dá o tamanho da matriz
no commit \texttt{80644ba9} do fork, contado em 2026-10-18.

\begin{longtable}[]{@{}
  >{\raggedright\arraybackslash}p{(\linewidth - 4\tabcolsep) * \real{0.7158}}
  >{\raggedright\arraybackslash}p{(\linewidth - 4\tabcolsep) * \real{0.2000}}
  >{\raggedright\arraybackslash}p{(\linewidth - 4\tabcolsep) * \real{0.0842}}@{}}
\caption{\label{tbl:regtests}Tamanho da matriz de regressão do WW3.
Fonte: ww3-gpu \texttt{regtests/README.md} em \texttt{58167d7},
contagens de \texttt{just\ ww3-gpu-matrix-list}. A contagem de todas as
entradas é o número de linhas que o \texttt{matrix.base} escreve, e uma
seção de máquina consegue desligar algumas dessas
linhas.}\tabularnewline
\toprule\noalign{}
\begin{minipage}[b]{\linewidth}\raggedright
Conjunto de entradas do \texttt{matrix.base}
\end{minipage} & \begin{minipage}[b]{\linewidth}\raggedright
Entradas
\end{minipage} & \begin{minipage}[b]{\linewidth}\raggedright
Fração
\end{minipage} \\
\midrule\noalign{}
\endfirsthead
\toprule\noalign{}
\begin{minipage}[b]{\linewidth}\raggedright
Conjunto de entradas do \texttt{matrix.base}
\end{minipage} & \begin{minipage}[b]{\linewidth}\raggedright
Entradas
\end{minipage} & \begin{minipage}[b]{\linewidth}\raggedright
Fração
\end{minipage} \\
\midrule\noalign{}
\endhead
\bottomrule\noalign{}
\endlastfoot
Todas as entradas & 758, em 60 testes & 100 \% \\
Entradas com a DIA (\texttt{NL1}) & 310 & 41 \% \\
Entradas com a física ST4 (\texttt{ST4}) & 136 & 18 \% \\
Entradas que o preset residente consegue selecionar & 7 (todas seriais)
& 0,9 \% \\
Regtests selecionados do ww3-gpu & 5, e o par \texttt{ww3\_ts1} & 0,7
\% \\
\end{longtable}

Uma execução das 310 entradas \texttt{NL1} leva dias. O conjunto
selecionado leva de uma a três horas numa estação de trabalho (Intel
i9-14900, 32 threads de hardware); esse tempo é uma estimativa, e não
uma medida ⚠. Por isso, uma execução da matriz completa não consegue
acompanhar cada commit. O ww3-gpu roda o conjunto selecionado como
proteção e adia a matriz completa para uma fase de validação. Até essa
fase, um argumento estático sustenta a paridade de bits: o GIMPLE em
\texttt{-O3} ou a desmontagem dos dois compiladores.

\section{O que este texto não afirma}\label{sec:limits}

\begin{enumerate}
\def\labelenumi{\arabic{enumi}.}
\tightlist
\item
  Nenhum dos dois projetos rodou o modelo inteiro numa GPU real. O PTX
  do chaceon nunca rodou. O ww3-gpu rodou um kernel de benchmark da DIA
  numa RTX 4090 (\citeproc{ref-diagpuv2}{Hoffmann 2026d}).
\item
  Este texto não afirma ganho de velocidade. As regras do ww3-gpu marcam
  os tempos registrados (\citeproc{ref-noaarecord}{Hoffmann 2026e}) como
  ``regtest pending'', e o chaceon não tem tempos medidos.
\item
  Numa GPU, a identidade de bits também exige que o \texttt{expf} e o
  \texttt{logf} do dispositivo devolvam os bits da glibc. O ww3-gpu
  manda os kernels para cópias de dispositivo das funções da glibc 2.44.
  Essas cópias são iguais às funções do host em toda entrada conferida.
  Nenhuma execução em dispositivo comparou as cópias ainda. Até uma
  execução comparar, uma tolerância fixa julga uma execução em GPU.
\item
  As duas referências são linhas diferentes do WW3 (NOAA-EMC e
  MetOcean). Por isso, a comparação entre os dois portes é por nome de
  rotina, e não por texto-fonte. Os dois registros contêm 85 rotinas.
\end{enumerate}

\section{Uso de IA}\label{sec:ai}

Agentes de programação Claude (Anthropic) produziram o código C++/Kokkos
e Rust, os shims, os scripts de comparação, as investigações e rascunhos
deste texto. Os agentes trabalharam sob a direção do primeiro autor, nos
dois repositórios. Pedro Cruz (GTA, UFRJ) é o orientador e Pedro Veras
Guimarães (LabECO, UFSC) é o coorientador do primeiro autor.

A maior parte dos commits veio do Claude Opus. Na experiência do
primeiro autor, o Claude Sonnet também portou uma quantidade grande de
código sem erros grosseiros ⚠. Um exemplo é a divisão do \texttt{W3SDS4}
em partes que rodam um ponto de mar por faixa vetorial (commit
\texttt{9cef132} do ww3-gpu, 841 linhas acrescentadas). Os arquivos de
saída desse commit foram idênticos byte a byte ao Fortran do mesmo
executável nos casos \texttt{smoke} e \texttt{full}.

Cada repositório publica as regras, as skills e os hooks dos agentes no
arquivo \texttt{AGENTS.md} e no diretório \texttt{.claude/}. Cada número
deste texto vem de um arquivo de um repositório, no commit que o texto
nomeia. Nada foi executado para este texto. Cada afirmação foi conferida
em 2026-11-11 contra a fonte e o commit que ela nomeia; um ⚠ marca uma
afirmação que não foi conferida numa fonte primária. A física do WW3 é
usada sem mudança; WAVEWATCH III é marca registrada da NOAA.

\textbf{Aviso.} Hoje, os dois repositórios estão configurados para o
Claude Code. Os autores vão continuar o trabalho para qualquer agente de
programação conseguir usar o ww3-gpu e o chaceon. Esse trabalho inclui
agentes que rodam modelos locais com o Ollama. O trabalho não está
completo, e este texto não afirma resultados para outros agentes nem
para modelos locais.

\section*{Referências}\label{referuxeancias}
\addcontentsline{toc}{section}{Referências}

\protect\phantomsection\label{refs}
\begin{CSLReferences}{1}{0}
\bibitem[\citeproctext]{ref-chaceon}
Hoffmann, Matheus. 2026a. {``{chaceon}: a {Rust} port of {WAVEWATCH
III}''}. \url{https://github.com/h0ffmann/chaceon}.

\bibitem[\citeproctext]{ref-oceanwaves}
---------. 2026b. {``Ocean waves: the physics of a spectral wave model
in {Lean} 4''}.
\url{https://github.com/h0ffmann/proof/tree/b863c15/ocean-waves}.

\bibitem[\citeproctext]{ref-ww3gpu}
---------. 2026c. {``{WW3 GPU Lab}: agentic coding and agentic research
on a bit-for-bit {GPU} port of {WAVEWATCH III}''}.
\url{https://doi.org/10.5281/zenodo.23221351}.

\bibitem[\citeproctext]{ref-diagpuv2}
---------. 2026d. {``{ww3-dia-gpu-v2}: the whole {W3SNL1}, inside
{WAVEWATCH III}, bit for bit''}.
\url{https://github.com/h0ffmann/ww3-dia-gpu-v2/tree/32aa2ef}.

\bibitem[\citeproctext]{ref-noaarecord}
---------. 2026e. {``A bit-identical {Kokkos} runtime for {WAVEWATCH
III}: {CPU} results and five speed-ups''}. ww3-gpu,
\texttt{docs/NOAA\_PROGRESS\_RECORD\_202610.md}, commit 58167d7.
\url{https://github.com/h0ffmann/ww3-gpu/blob/58167d7/docs/NOAA_PROGRESS_RECORD_202610.md}.

\bibitem[\citeproctext]{ref-adr0001}
---------. 2026f. {``{ADR-0001}: no proof language for the {GPU} port;
{SMT} queries and exhaustive sweeps instead''}. ww3-gpu,
\texttt{docs/ADRs/ADR-0001-proof-language.md}, commit 58167d7.
\url{https://github.com/h0ffmann/ww3-gpu/blob/58167d7/docs/ADRs/ADR-0001-proof-language.md}.

\bibitem[\citeproctext]{ref-koldunov2026}
Koldunov, Nikolay V., Suvarchal K. Cheedela, Sergey Danilov, Dmitry
Sidorenko, Sebastian Beyer, e Thomas Jung. 2026. {``An Ocean Model
Ported by a Large Language Model: Experience and Lessons from {FESOM2}
({Fortran} to {C} to {C++/Kokkos})''}. arXiv:2606.11356
{[}physics.ao-ph{]}. \url{https://doi.org/10.48550/arXiv.2606.11356}.

\bibitem[\citeproctext]{ref-renomo}
LabECO/UFSC. 2025. {``{ReNOMO}: Rede Nacional de Observação e
Monitoramento Oceânico''}. \url{https://labeco.ufsc.br/renomo/}.

\bibitem[\citeproctext]{ref-ww4}
NOAA Environmental Modeling Center. 2026. {``{WAVEWATCH IV}''}.
\url{https://github.com/NOAA-EMC/WW4}.

\bibitem[\citeproctext]{ref-on525}
Tolman, Hendrik L. 2025. {``The {WAVEWATCH III} Software Modernization
Project: Phase {I} report''}. Office Note 525. College Park, MD:
NOAA/NWS/NCEP. \url{https://doi.org/10.25923/h7j3-1h25}.

\bibitem[\citeproctext]{ref-ondasventos}
Universidade Federal de Santa Catarina. 2026. {``Ondas \& Ventos''}.
\url{https://ondaseventos.ufsc.br/}.

\end{CSLReferences}

\end{document}
